\chapter{Getting Access: FX}\label{ch:m2:getting-access-fx}

On 15 January 2015 the Swiss National Bank stopped capping the franc, and the franc
jumped (chapter 18). By the BIS's account, the moves jolted the prime-brokerage
industry, the banks that stand behind their clients' trades, and the prime brokers
reassessed the business: they kept their largest clients, the big
principal trading firms, and shed retail aggregators, smaller hedge funds and some
high-frequency traders, raising capital requirements, tightening admission and raising
fees; FX volume through prime brokers fell by 22\% between 2013 and 2016. A new fund
whose strategy trades a billion dollars a day must first find someone to stand behind
its trades, and then live within the limits that person sets. This chapter explains
how: prime-broker credit and \omterm{def:m2:getting-access-fx:giveup}{give-up lines}, the limits and the gate that enforces them,
becoming a liquidity provider, and the disclosures and reviews on both sides.

\section{Prime-broker credit and give-up lines}

A fund cannot trade FX with dozens of dealers, each of which would have to assess its
credit, sign documentation and set a limit. An FX prime broker (chapter 14) does it
once: the fund trades with any dealer that accepts the prime broker's name, and each
trade is \emph{given up} to the prime broker, which becomes the dealer's counterparty
and books an offsetting trade with the fund.

\begin{definition}[Give-up line, designation notice]\label{def:m2:getting-access-fx:giveup}
A \emph{give-up line}\index{give-up line} is a prime broker's agreement to accept the
trades that a client executes with a given executing dealer, within limits. A
\emph{designation notice}\index{designation notice} is the prime broker's notice to an
executing dealer that designates the client as a party authorised to trade on its
behalf, and sets the trade types, maximum tenors, currencies and limits of that
authority.
\end{definition}

The standard documents were published by the Foreign Exchange Committee in New York in
2005: in the Master FX Give-Up Agreement, the prime broker authorises each designated
party, by a notice in a standard form, to trade with the dealer on its behalf; the
authority is limited to the types, tenors, currencies and offices in the notice and to
two limits, and the prime broker is not liable for a trade that takes a limit above
its level, or further above it, without its consent
(\cref{fig:m2:getting-access-fx:giveup}). The client's cost is the prime broker's fees,
per million traded and for the credit it extends, and the prime broker's right to cut
the limits when it wants.

\begin{omfigure}
\centering
\begin{tikzpicture}[font=\small,
  box/.style={draw,thick,rounded corners=2pt,align=center,minimum height=1cm,minimum width=2.8cm,inner sep=3pt},
  arr/.style={-{Stealth[length=2.5mm]},thick}]
\node[box,draw=omDef,fill=omDef!8] (fund) at (0,0) {fund\\(designated party)};
\node[box,draw=gray,fill=gray!8] (ed) at (5.4,1.4) {executing dealer 1};
\node[box,draw=gray,fill=gray!8] (ed2) at (5.4,-1.4) {executing dealer 2};
\node[box,draw=omThm,fill=omThm!8] (pb) at (10.8,0) {prime broker};
\draw[arr] (fund) -- node[above,sloped,font=\footnotesize] {executes} (ed);
\draw[arr] (fund) -- node[below,sloped,font=\footnotesize] {executes} (ed2);
\draw[arr] (ed) -- node[above,sloped,font=\footnotesize] {gives up} (pb);
\draw[arr] (ed2) -- node[below,sloped,font=\footnotesize] {gives up} (pb);
\draw[arr,dashed,omThm] (pb.north) |- node[pos=0.75,above,font=\footnotesize] {designation notice: pairs, tenors, NOP and settlement limits} (ed.north);
\draw[arr,omThm] (pb.south) |- node[pos=0.75,below,font=\footnotesize] {offsetting trade with the fund} ([yshift=-2.4cm]fund.south) -- (fund.south);
\end{tikzpicture}

\omcaption{A \omterm{def:m2:getting-access-fx:giveup}{give-up line}. The fund executes with any dealer that accepts its prime
broker's name; each trade is given up to the prime broker, which faces the dealer and
books an offsetting trade with the fund. The prime broker's notice tells each dealer
what the fund may trade on its behalf and within which limits. Schematic.}\label{fig:m2:getting-access-fx:giveup}
\end{omfigure}

\begin{definition}[Prime-of-prime]\label{def:m2:getting-access-fx:pop}
\emph{Prime-of-prime}\index{prime-of-prime} is an arrangement in which a client is
prime-brokered by a non-dealer bank or firm that is itself prime-brokered by an FX
dealing bank.
\end{definition}

After 2015 the clients the banks shed turned to \omterm{def:m2:getting-access-fx:pop}{prime-of-prime} providers, which take
their credit and bundle it into their own relationship with a bank
(\cref{fig:m2:getting-access-fx:pop}). It adds a layer, and its costs and limits, but
it is often the only way in for a small firm. The New York committee published in 2014
sample certifications and best practices for such intermediated arrangements, in
which the intermediary certifies its relationships to the prime brokers and dealers.

\begin{omfigure}
\centering
\begin{tikzpicture}[font=\small,
  box/.style={draw,thick,rounded corners=2pt,align=center,minimum height=1cm,minimum width=2.6cm,inner sep=3pt},
  arr/.style={-{Stealth[length=2.5mm]},thick}]
\node[box,draw=omDef,fill=omDef!8] (c) at (0,0) {small fund or\\retail aggregator};
\node[box,draw=omProp,fill=omProp!8] (pop) at (4.6,0) {prime-of-prime\\(non-dealer)};
\node[box,draw=omThm,fill=omThm!8] (pb) at (9.2,0) {FX dealing bank\\(prime broker)};
\node[box,draw=gray,fill=gray!8] (v) at (13.4,0) {venues and\\dealers};
\draw[arr] (c) -- node[above,font=\footnotesize] {credit} (pop);
\draw[arr] (pop) -- node[above,font=\footnotesize] {credit} (pb);
\draw[arr] (pb) -- node[above,font=\footnotesize] {name} (v);
\end{tikzpicture}

\omcaption{\omterm{def:m2:getting-access-fx:pop}{Prime-of-prime}: a client too small for a bank's prime brokerage is
prime-brokered by a non-dealer firm, which is itself prime-brokered by a dealing bank
and reaches venues and dealers under its name. Schematic.}\label{fig:m2:getting-access-fx:pop}
\end{omfigure}

\section{Limits: net open position and settlement}

\begin{definition}[Net open position limit, settlement limit]\label{def:m2:getting-access-fx:limits}
A \emph{net open position limit}\index{net open position limit} caps the client's net
currency exposure through the prime broker, measured in one currency across all
currencies. A \emph{settlement limit}\index{settlement limit} caps, for each value date,
the amount the prime broker will have to settle on the client's trades.
\end{definition}

The Master Agreement defines both in dollar values: the net open position aggregates,
by a method set in the schedule, the net dollar value of each currency across the
designated party's trades; the net daily settlement amount of a value date is the sum
of the dollar values of the currencies with a net amount owed to the prime broker that
day. This chapter's gate takes the net open position as the sum of the net long dollar
values, which equals the sum of the net shorts when every currency, the dollar
included, is counted, and the settlement amount as the sum of the currencies to be
received on the value date.

\begin{example}[Two trades through the gate]\label{ex:m2:getting-access-fx:nop}
Prime broker A allows EURUSD and USDJPY up to seven days, a net open position of USD~100
million and settlement of USD~500 million per value date. The fund buys EUR~50 million
at 1.10 for spot: it is long EUR worth USD~55 million and short USD~55 million, a net
open position of 55 million and a settlement amount of 55 million. It then buys USD~20
million against yen: the dollar short shrinks to 35 million and a yen short of 20
million appears; the longs are still the 55 million of euros, and so is the settlement
amount. A further purchase of EUR~95 million would take the position to 159.5 million
and is rejected; a sale of euros would be accepted even if the limit had been cut below
the position, because it reduces it.
\end{example}

A limit is only as good as its enforcement. The fund's own gate must check every order
before it leaves, in the order path, against what each prime broker's notice allows:
a dealer that accepts a trade beyond the limit may find the prime broker refusing it,
leaving the fund facing the dealer directly or not at all. The build of this chapter
is that gate, written in C++20 and Rust for the order path and in Python for analysis
(\cref{fig:m2:getting-access-fx:path}).

\begin{omfigure}
\begin{tikzpicture}
\begin{axis}[width=11cm,height=5cm,
  xlabel={orders in the day},ylabel={NOP used (\% of limit)},
  tick label style={font=\small},label style={font=\small},
  xmin=0,xmax=240,ymin=0,ymax=105,grid=major,grid style={gray!25},
  legend columns=3,legend style={font=\footnotesize,at={(0.5,-0.3)},anchor=north,draw=none,fill=none,/tikz/every even column/.append style={column sep=8pt}},
  table/col sep=comma]
\addplot[omDef,thick] table[x=order,y=A]{figdata/markets-2/27-getting-access-fx/nop_path.csv};
\addplot[omThm,thick] table[x=order,y=B]{figdata/markets-2/27-getting-access-fx/nop_path.csv};
\addplot[omProp,thick] table[x=order,y=C]{figdata/markets-2/27-getting-access-fx/nop_path.csv};
\legend{A (fee 3),B (fee 5),C (fee 8)}
\end{axis}
\end{tikzpicture}

\omcaption{A simulated day of 240 orders routed by the gate to the cheapest prime broker
whose limits allow them: the net open position at each, in per cent of its limit. A
fills first, from order 107 the flow that adds to the position spills to B, and then to
C, whose position ends the day at 76\% of its limit. No order is rejected.
Illustrative; data: the chapter's tutorial, seeded.}\label{fig:m2:getting-access-fx:path}
\end{omfigure}

\section{Becoming a liquidity provider}

A fund that quotes, rather than only takes prices, must also be accepted as a
liquidity provider: on a venue, which lets it stream prices into the venue's book; or
directly to clients, through a \omterm{def:m2:the-fx-market:sdp}{single-dealer platform} or an aggregator.

\begin{definition}[Liquidity-provider review]\label{def:m2:getting-access-fx:review}
A \emph{liquidity-provider review}\index{liquidity-provider review} is the periodic
assessment, by a venue or a client, of a liquidity provider's quoting: how often its
prices are hit and filled, its \omterm{def:m2:fx-spot-microstructure:lastlook}{reject rate} and hold times under \omterm{def:m2:fx-spot-microstructure:lastlook}{last look}, and the
mark-outs of the trades done with it.
\end{definition}

The metrics are those of chapter 15, seen from the other side. A liquidity provider
that rejects too often, holds requests too long, or whose fills are systematically
followed by moves against its clients may be ranked down in aggregators' routing or
removed from a venue's stream; one that shows firm, well-used prices can expect more
flow. It needs its own credit to do this: every trade it makes as a maker is still
given up to its prime broker and counts against the same limits.

\section{Brokerage, disclosures and reviews}

The \omterm{def:m2:fx-spot-microstructure:code}{FX Global Code} asks market participants to disclose how they trade, and in August
2021 the Global Foreign Exchange Committee published standard disclosure cover sheets
for liquidity providers and for trading platforms, with a guidance paper on \omterm{def:m2:fx-spot-microstructure:lastlook}{last look}.
A liquidity provider's cover sheet states in one place whether it acts as principal or
agent, whether it shares data from client interactions, whether it pre-hedges, whether
and how it uses \omterm{def:m2:fx-spot-microstructure:lastlook}{last look} (symmetrically or not, the maximum and minimum window, and
whether it trades during the window), and where to find its policies on aggregation,
discretion, time-stamping, stop-loss orders, partial fills, reference prices, mark-ups
and the internal sharing of confidential information. A fund reads these before it
chooses whom to trade with, and must write its own if it becomes a liquidity provider.

\begin{dated}{2026-09}{FX access documents}\label{dat:m2:getting-access-fx:disclosure}
Master FX Give-Up Agreement: published by the Foreign Exchange Committee in 2005, with a
standard notice form. \omterm{def:m2:fx-spot-microstructure:code}{FX Global Code}: last updated in December 2024, with a December
2024 liquidity-provider disclosure cover sheet; cover sheets first published in August
2021.
\end{dated}

\section{Tutorial: a pre-trade limit gate}\label{tut:m2:getting-access-fx:gate}

\begin{tutorial}
\textbf{Goal.} Implement a pre-trade check that enforces each prime broker's allowed
pairs, tenors, net open position and \omterm{def:m2:getting-access-fx:limits}{settlement limits}, route each order to the
cheapest prime broker that accepts it, and plan the allocation of a day's flow.
\textbf{End state:} \cref{fig:m2:getting-access-fx:path,fig:m2:getting-access-fx:plans},
\cref{ex:m2:getting-access-fx:nop} and the numbers of the weekend problem.
\begin{tutsteps}
\item \textbf{The gate}, in C++20: the check, which leaves the book unchanged, and the
update once the trade is done.
\omcode{code/firm/pblimits/cpp/firm_pblimits.hpp}{46}{75}{The prime-broker book and its pre-trade check, C++20.}\label{lst:m2:getting-access-fx:gate}
\item \textbf{The router} and the limits file, in Python.
\omcode{code/firm/pblimits/firm_pblimits.py}{85}{102}{Routing to the cheapest prime broker with room, and reading the limits.}\label{lst:m2:getting-access-fx:route}
\item \textbf{Run} \texttt{access\_demo.plan()}, \texttt{access\_demo.brute\_force()},
\texttt{access\_demo.day\_of\_orders()} and \texttt{fig\_access.py}; build and test the
C++ and Rust gates.
\end{tutsteps}
\textbf{What to change next.} Add a maximum order size per pair to each notice; then
make the router prefer, among the prime brokers that accept an order, the one whose
overnight charge is lowest when the order adds to the position, and measure the cost
saved.
\end{tutorial}

\section{Build: the prime-broker limit gate}\label{bld:m2:getting-access-fx:pblimits}

\begin{build}
\bfield{Purpose} Every order the miniature firm sends in FX passes this gate before it
leaves: it must fit a prime broker's notice, and is given up to the cheapest one that
accepts it.
\bfield{Interface} C++20 \texttt{firm::pblimits::Book} with \texttt{check(order,
usd\_per)}, \texttt{apply}, \texttt{nop()}, \texttt{settlement(day)}, and
\texttt{route(books, order, usd\_per)}; the same in Rust; Python \texttt{Book},
\texttt{check}, \texttt{apply}, \texttt{route} and \texttt{read\_limits(path)} for the
shared limits file \texttt{code/firm/pblimits/limits.csv}.
\bfield{Rules} Dollar values at given rates; net open position the sum of net longs;
settlement the sum of net receipts per value date; reject if a pair or tenor is not
allowed, or if a limit is exceeded or further exceeded; route by increasing fee.
\bfield{Acceptance tests} The same in the three languages: the net open position and
settlement of two trades; rejection above the limit and acceptance of a reducing trade
after a cut; the \omterm{def:m2:getting-access-fx:limits}{settlement limit} per value date; pair and tenor rejections; routing
to the cheapest broker with room.
\bfield{Stretch} Limits per executing dealer as well as per prime broker; a
lock-free book for concurrent order threads (One Quant Book 13); end-of-day moves of
positions between prime brokers; limits on options by their delta and premium.
\end{build}

\begin{omsources}
\item Foreign Exchange Committee, Master FX Give-Up Agreement (2005); FXC and FMLG,
market practice memoranda of May 2013 and July 2014 on prime-broker notices and
intermediated arrangements.
\item M.~Moore, A.~Schrimpf and V.~Sushko, ``Downsized FX markets: causes and
implications'', BIS Quarterly Review, December 2016.
\item Global Foreign Exchange Committee, press release of 18 August 2021 and disclosure
cover sheets; a published liquidity-provider cover sheet.
\end{omsources}

\section{Exercises}

\begin{exercise}[$\star$]\label{exo:m2:getting-access-fx:1}
A fund's book is long EUR worth USD~40 million, short USD~25 million and short JPY worth
USD~15 million. What is its net open position, and why do the longs equal the shorts?
\end{exercise}

\begin{exercise}[$\star$]\label{exo:m2:getting-access-fx:2}
What does a \omterm{def:m2:getting-access-fx:giveup}{designation notice} tell an executing dealer?
\end{exercise}

\begin{exercise}[$\star$]\label{exo:m2:getting-access-fx:3}
Why did prime brokers shed smaller clients after 2015, and where did those clients go?
\end{exercise}

\begin{exercise}[$\star\star$]\label{exo:m2:getting-access-fx:4}
For the same value date the fund buys EUR~100 million and sells EUR~60 million against
dollars at 1.10. What is the settlement amount, and what would it be if the two trades
were at different prime brokers?
\end{exercise}

\begin{exercise}[$\star\star$]\label{exo:m2:getting-access-fx:5}
Why does the gate accept a trade that reduces a position already above its limit?
\end{exercise}

\begin{exercise}[$\star\star$]\label{exo:m2:getting-access-fx:6}
Which numbers would a venue look at to review the fund as a liquidity provider?
\end{exercise}

\begin{exercise}[$\star\star\star$]\label{exo:m2:getting-access-fx:7}
\emph{Coding.} Cut prime broker A's \omterm{def:m2:getting-access-fx:limits}{net open position limit} to USD~50 million and run the
day of orders again. How much more flow goes to C, and where does the day's position
end?
\end{exercise}

\begin{exercise}[$\star\star\star$]\label{exo:m2:getting-access-fx:8}
\emph{Find the flaw.} ``Our strategy trades only liquid majors; one prime broker is
enough.'' Correct it.
\end{exercise}

\section{Problem: Three Prime Brokers}

\begin{problem}[{Weekend problem --- keeping every limit at least cost}]\label{pb:m2:getting-access-fx:1}
A fund's day, in USD millions: 1\,000 of flow in EURUSD and USDJPY, 150 in GBPUSD and 50
in USDMXN, and an overnight net open position of 300, of which 30 in pesos. Three prime
brokers' notices, from \texttt{limits.csv}: A charges 3 dollars per million traded,
allows EURUSD and USDJPY up to 7 days, a net open position of 100 and settlement of 500;
B charges 5, adds GBPUSD, up to 30 days, 150 and 600; C charges 8, adds USDMXN, up to a
year, 300 and 1\,000. Overnight positions cost 30, 20 and 50 dollars per million per day
at A, B and C. For planning, take the settlement amount at a broker to be the whole flow
given up there, without netting.

\medskip\noindent\textbf{Part I --- The notices.}
\begin{enumerate}
\item Which brokers can take each part of the flow?
\item Compute the net open position after buying EUR~50 million at 1.10 and then USD~20
million against yen.
\item At B, after buying EUR~500 million for spot, can the fund buy 50 million more for
the same value date? For the next day?
\item A's limit is cut to 50 million while the fund is long 99 million: which orders
still pass?
\item Why does the notice limit tenors?
\end{enumerate}

\medskip\noindent\textbf{Part II --- The plan.}
\begin{enumerate}[resume]
\item Allocate the flow at least cost within the \omterm{def:m2:getting-access-fx:limits}{settlement limits}.
\item Allocate the overnight position at least cost within the \omterm{def:m2:getting-access-fx:limits}{net open position limits}.
\item What does the plan cost per day, and per year of 250 days?
\item What would it cost to do everything at C, the only broker that takes every
currency?
\item Check the plan against a search over all splits in steps of 10 million.
\end{enumerate}

\medskip\noindent\textbf{Part III --- The day.}
\begin{enumerate}[resume]
\item Route a simulated day of 240 orders to the cheapest broker with room: when does A's
position first reach its limit?
\item How many orders are rejected?
\item Where does the day's position end, and why at the most expensive broker?
\item Why can the router's settlement amounts be far below the flow it sends to a broker?
\item How would you bring the router closer to the plan?
\end{enumerate}

\medskip\noindent\textbf{Part IV --- Judgement.}
\begin{enumerate}[resume]
\item What happens to the fund if B cuts its limits by half after a shock?
\item Why is a \omterm{def:m2:getting-access-fx:pop}{prime-of-prime} arrangement an alternative, and what does it add?
\item What would you disclose, and read, before becoming a liquidity provider?
\item State the \emph{named result}: the allocation of flow and position that keeps every
limit at least cost, and that cost.
\item In one sentence: what does a fund need before its first FX trade?
\end{enumerate}
\end{problem}

\begin{omfigure}
\begin{tikzpicture}
\begin{axis}[width=11cm,height=4.8cm,ybar stacked,bar width=26pt,
  ylabel={cost per day (USD thousand)},
  tick label style={font=\small},label style={font=\small},
  xtick=data,xticklabels={all at C,{fees\\optimised},{NOP\\optimised},{both\\optimised}},xticklabel style={align=center},
  ymin=0,ymax=27,grid=major,grid style={gray!25},enlarge x limits=0.18,
  legend columns=2,legend style={font=\footnotesize,at={(0.5,-0.38)},anchor=north,draw=none,fill=none,/tikz/every even column/.append style={column sep=8pt}},
  table/col sep=comma]
\addplot[fill=omDef!60,draw=omDef] table[x=k,y=fees]{figdata/markets-2/27-getting-access-fx/plans.csv};
\addplot[fill=omThm!50,draw=omThm] table[x=k,y=nop]{figdata/markets-2/27-getting-access-fx/plans.csv};
\legend{fees on flow,overnight position charges}
\end{axis}
\end{tikzpicture}

\omcaption{Daily cost of the weekend problem's fund under four plans: everything at C;
the flow allocated at least fee cost with the position at C; the position allocated at
least charge with the flow at C; and both allocated. The best plan costs USD~13\,800 a
day against 24\,600. Illustrative; data: the chapter's tutorial.}\label{fig:m2:getting-access-fx:plans}
\end{omfigure}

\section{Interview questions}

\begin{interviewq}[$\star$ \iqroles{trader, developer}]\label{iq:m2:getting-access-fx:1}
What does an FX prime broker do for a fund, and what does a give-up mean?
\end{interviewq}

\begin{interviewq}[$\star$ \iqroles{risk, developer}]\label{iq:m2:getting-access-fx:2}
What are a \omterm{def:m2:getting-access-fx:limits}{net open position limit} and a \omterm{def:m2:getting-access-fx:limits}{settlement limit}, and how do you compute them?
\end{interviewq}

\begin{interviewq}[$\star\star$ \iqroles{developer}]\label{iq:m2:getting-access-fx:3}
Where in the order path should the limit check sit, and what must it do when it cannot
decide in time?
\end{interviewq}

\begin{interviewq}[$\star\star$ \iqroles{trader, risk}]\label{iq:m2:getting-access-fx:4}
How would you allocate a fund's flow across several prime brokers?
\end{interviewq}

\begin{interviewq}[$\star\star$ \iqroles{trader}]\label{iq:m2:getting-access-fx:5}
What should a fund ask a liquidity provider before trading with it?
\end{interviewq}

\begin{interviewq}[$\star\star\star$ \iqroles{developer}]\label{iq:m2:getting-access-fx:6}
Design a limit gate that checks 100\,000 orders a second across ten prime brokers with
a worst-case latency of five microseconds.
\end{interviewq}
